\newcommand{\figuraMasaPlanetaDosCapas}{%
\begin{figure}[H]
    \centering
    \begin{tikzpicture}[>=Latex,font=\small,line cap=round,line join=round]
        \coordinate (O) at (3.05,3.05);

        \fill[orange!16] (O) circle (2.45);
        \fill[green!18] (O) circle (1.72);
        \fill[blue!22] (O) circle (0.95);

        \draw[very thick] (O) circle (2.45);
        \draw[very thick,dashed,green!45!black] (O) circle (1.72);
        \draw[very thick,blue!65!black] (O) circle (0.95);
        \fill[black] (O) circle (0.06);
        \node[anchor=north west] at ($(O)+(0.05,-0.05)$) {$O$};

        \node[blue!65!black] at (3.05,3.55) {$\rho_1$};
        \node[green!35!black] at (2.60,1.75) {$\rho_2$};
        \node[orange!70!black] at (2.60,5.00) {$\rho_2$};

        \draw[->,thick] (O)--++(0:2.68);
        \draw[thick,blue!65!black]
            ($(O)+(0:0.95)+(0,-0.10)$)--++(0,0.20);
        \node[anchor=north,blue!65!black]
            at ($(O)+(0:0.95)+(0,-0.14)$) {$a_1$};

        \draw[thick,green!45!black]
            ($(O)+(0:1.72)+(0,-0.10)$)--++(0,0.20);
        \node[anchor=north,green!35!black]
            at ($(O)+(0:1.72)+(0,-0.14)$) {$r$};

        \draw[thick]
            ($(O)+(0:2.45)+(0,-0.10)$)--++(0,0.20);
        \node[anchor=north]
            at ($(O)+(0:2.45)+(0,-0.14)$) {$a$};

        \node[font=\bfseries] at (9.05,5.55) {Regiones del planeta};

        \fill[blue!22] (6.55,4.92) rectangle (6.95,5.22);
        \draw[blue!65!black] (6.55,4.92) rectangle (6.95,5.22);
        \node[anchor=west] at (7.12,5.07)
            {núcleo completo: $0<r'<a_1$};

        \fill[green!18] (6.55,4.30) rectangle (6.95,4.60);
        \draw[green!45!black,dashed] (6.55,4.30) rectangle (6.95,4.60);
        \node[anchor=west] at (7.12,4.45)
            {capa encerrada: $a_1<r'<r$};

        \fill[orange!16] (6.55,3.68) rectangle (6.95,3.98);
        \draw[orange!70!black] (6.55,3.68) rectangle (6.95,3.98);
        \node[anchor=west] at (7.12,3.83)
            {región exterior: $r<r'<a$};

        \node[font=\bfseries] at (9.12,2.88)
            {Masa encerrada para $a_1<r<a$};
        \node[align=center] at (9.12,1.95)
            {$\displaystyle
            M(r)=
            \underbrace{\frac{4\pi}{3}\rho_1a_1^3}_{
            \textcolor{blue!65!black}{\text{núcleo}}}
            +
            \underbrace{\frac{4\pi}{3}\rho_2(r^3-a_1^3)}_{
            \textcolor{green!35!black}{\text{capa hasta }r}}$};
        \node[text=black!65] at (9.12,1.08)
            {La región naranja no se suma.};
    \end{tikzpicture}
    \caption{La circunferencia discontinua representa la esfera gaussiana
    de radio $r$. Para $a_1<r<a$, se suman la región azul y la región verde;
    la región naranja queda fuera y no se incluye en $M(r)$.}
    \label{fig:masa-planeta-dos-capas}
\end{figure}%
}